% Grammatikerklaerung nach der hochgeladenen Erklaerungsvorlage

\documentclass[a4paper,11pt]{article}

\usepackage[a4paper,margin=2cm]{geometry}
\usepackage[T1]{fontenc}
\usepackage[utf8]{inputenc}
\usepackage[ngerman,english]{babel}
\usepackage{lmodern}
\usepackage{microtype}
\usepackage[table]{xcolor}
\usepackage{parskip}
\usepackage{enumitem}
\usepackage{array}
\usepackage{longtable}
\usepackage[most]{tcolorbox}
\usepackage{titlesec}
\usepackage[hidelinks]{hyperref}

\definecolor{HeaderBlueDark}{RGB}{70,110,170}
\definecolor{RowBlueLight}{RGB}{235,243,255}
\definecolor{RowBlueDark}{RGB}{220,233,250}
\definecolor{SoftGray}{RGB}{90,90,90}

\setlength{\parindent}{0pt}
\setlist[itemize]{leftmargin=1.4em,itemsep=0.35em,topsep=0.35em}
\renewcommand{\arraystretch}{1.35}
\setlength{\tabcolsep}{5pt}
\linespread{1.06}

\titleformat{\section}[block]
{\Large\bfseries\color{HeaderBlueDark}}
{}{0pt}{}
\titlespacing{\section}{0pt}{1.3em}{0.8em}

\titleformat{\subsection}[block]
{\large\bfseries\color{HeaderBlueDark}}
{}{0pt}{}
\titlespacing{\subsection}{0pt}{1.0em}{0.6em}

\newtcolorbox{exbox}{enhanced,breakable,colback=RowBlueLight,colframe=RowBlueDark,boxrule=0.4pt,arc=1.5mm,left=1.2mm,right=1.2mm,top=1mm,bottom=1mm,title={Beispiele \textnormal{(Examples)}},fonttitle=\bfseries,colbacktitle=RowBlueDark,coltitle=black}

\newtcolorbox{warningbox}{enhanced,colback=RowBlueLight,colframe=HeaderBlueDark,boxrule=0.6pt,arc=1.5mm,left=1.5mm,right=1.5mm,top=1mm,bottom=1mm,title={Achtung \textnormal{(Important)}},fonttitle=\bfseries,colbacktitle=HeaderBlueDark,coltitle=white}

\NewDocumentEnvironment{grammarentry}{mm+b}{%
    \begin{tcolorbox}[enhanced,breakable,colback=white,colframe=HeaderBlueDark,boxrule=0.6pt,arc=1.5mm,left=1.5mm,right=1.5mm,top=1mm,bottom=1mm,colbacktitle=HeaderBlueDark,coltitle=white,fonttitle=\bfseries,title={#1\hfill\normalfont\footnotesize #2}]
    #3
    \end{tcolorbox}
}{}

\newcommand{\GermanText}[1]{\textbf{Deutsch: }#1\par}
\newcommand{\EnglishText}[1]{{\small\color{SoftGray}\textbf{English: }#1\par}}
\newcommand{\ExampleItem}[2]{\item \textbf{#1}\\{\small\color{SoftGray}#2}}

\title{Die Grammatik von \glqq vor\grqq}
\author{}
\date{}

\begin{document}
    \maketitle
    \tableofcontents
    \newpage

    \section{Grundbedeutungen}

        \begin{grammarentry}{vor}{Präposition, Adverb und Verbpräfix}
            \GermanText{Das Wort \textbf{vor} hat mehrere Funktionen. Am wichtigsten ist \textbf{vor} als Präposition. Es kann räumlich \glqq in front of\grqq, zeitlich \glqq before\grqq{} oder \glqq ago\grqq{} und kausal etwa \glqq because of\grqq{} bedeuten.}
            \EnglishText{The word \textbf{vor} has several functions. Most importantly, \textbf{vor} is a preposition. It can mean ``in front of'' spatially, ``before'' or ``ago'' temporally, and approximately ``because of'' when expressing a cause.}

            \GermanText{Außerdem erscheint \textbf{vor} als Adverb und als Bestandteil vieler trennbarer Verben, zum Beispiel \textbf{vorlesen}, \textbf{vorstellen} und \textbf{vorbereiten}.}
            \EnglishText{It also appears as an adverb and as a component of many separable verbs, such as \textbf{vorlesen}, \textbf{vorstellen}, and \textbf{vorbereiten}.}
        \end{grammarentry}

        \begin{exbox}
            \begin{itemize}
                \ExampleItem{Das Auto steht vor dem Haus.}{The car is in front of the house.}
                \ExampleItem{Wir sprechen vor dem Essen.}{We speak before the meal.}
                \ExampleItem{Ich war vor zwei Jahren in Berlin.}{I was in Berlin two years ago.}
                \ExampleItem{Er zittert vor Angst.}{He is trembling with fear.}
            \end{itemize}
        \end{exbox}


    \section{Räumliches \glqq vor\grqq: Dativ oder Akkusativ}

        \begin{grammarentry}{Wechselpräposition}{Wo? = Dativ; Wohin? = Akkusativ}
            \GermanText{Bei einer räumlichen Bedeutung ist \textbf{vor} eine Wechselpräposition. Der Kasus hängt nicht einfach vom Verb ab, sondern davon, ob der Satz eine \textbf{Position} oder eine \textbf{gerichtete Veränderung der Position} beschreibt.}
            \EnglishText{With a spatial meaning, \textbf{vor} is a two-way preposition. The case does not simply depend on the verb, but on whether the sentence describes a \textbf{position} or a \textbf{directed change of position}.}

            \GermanText{Frage \textbf{Wo?}: Dativ. Etwas oder jemand befindet sich bereits an einem Ort vor etwas anderem.}
            \EnglishText{Question \textbf{Where?}: dative. Something or someone is already located in a position in front of something else.}

            \GermanText{Frage \textbf{Wohin?}: Akkusativ. Etwas oder jemand wird an einen Ort vor etwas anderem gebracht oder bewegt.}
            \EnglishText{Question \textbf{Where to?}: accusative. Something or someone is brought or moved to a position in front of something else.}
        \end{grammarentry}

        \rowcolors{2}{RowBlueLight}{RowBlueDark}
        \begin{longtable}{>{\bfseries}p{2.5cm}p{3cm}p{4.8cm}p{4.1cm}}
            \rowcolor{HeaderBlueDark}
            \color{white}\textbf{Bedeutung} &
            \color{white}\textbf{Frage / Kasus} &
            \color{white}\textbf{Beispiel} &
            \color{white}\textbf{English} \\
            Position & Wo? + Dativ & Der Bus wartet vor dem Bahnhof. & The bus is waiting in front of the station. \\
            Richtung / Platzierung & Wohin? + Akkusativ & Der Fahrer stellt den Bus vor den Bahnhof. & The driver positions the bus in front of the station. \\
        \end{longtable}

        \begin{exbox}
            \begin{itemize}
                \ExampleItem{Ich stehe vor der Tür.}{I am standing in front of the door.}
                \ExampleItem{Ich stelle mich vor die Tür.}{I position myself in front of the door.}
                \ExampleItem{Das Fahrrad steht vor dem Gebäude.}{The bicycle is in front of the building.}
                \ExampleItem{Ich stelle das Fahrrad vor das Gebäude.}{I put the bicycle in front of the building.}
            \end{itemize}
        \end{exbox}

        \begin{warningbox}
            \GermanText{Bewegung bedeutet nicht automatisch Akkusativ. Entscheidend ist, ob sich die Position relativ zum Bezugsobjekt ändert. In \glqq Ich gehe vor dem Haus spazieren\grqq{} bleibt die Handlung im Bereich vor dem Haus; deshalb steht der Dativ.}
            \EnglishText{Movement does not automatically require the accusative. What matters is whether the position relative to the reference object changes. In ``Ich gehe vor dem Haus spazieren,'' the action remains in the area in front of the house; therefore, the dative is used.}
        \end{warningbox}


    \section{Artikel und Formen nach \glqq vor\grqq}

        \begin{grammarentry}{Kasusformen}{bestimmter Artikel}
            \GermanText{Bei der räumlichen Verwendung muss man die Dativ- und Akkusativformen der Artikel unterscheiden. Im Plural erhält das Nomen im Dativ normalerweise zusätzlich die Endung \textbf{-n}, sofern es sie nicht schon hat.}
            \EnglishText{With spatial usage, one must distinguish the dative and accusative forms of the articles. In the dative plural, the noun normally also receives the ending \textbf{-n}, unless it already has it.}
        \end{grammarentry}

        \rowcolors{2}{RowBlueLight}{RowBlueDark}
        \begin{longtable}{>{\bfseries}p{2.6cm}p{2.8cm}p{2.8cm}p{2.8cm}p{3.2cm}}
            \rowcolor{HeaderBlueDark}
            \color{white}\textbf{Kasus} &
            \color{white}\textbf{Maskulin} &
            \color{white}\textbf{Feminin} &
            \color{white}\textbf{Neutrum} &
            \color{white}\textbf{Plural} \\
            Dativ & vor dem Mann & vor der Frau & vor dem Kind & vor den Häusern \\
            Akkusativ & vor den Mann & vor die Frau & vor das Kind & vor die Häuser \\
        \end{longtable}

        \begin{exbox}
            \begin{itemize}
                \ExampleItem{Der Hund sitzt vor dem Mann.}{The dog is sitting in front of the man.}
                \ExampleItem{Der Hund läuft vor den Mann.}{The dog runs to a position in front of the man.}
                \ExampleItem{Wir warten vor den Häusern.}{We are waiting in front of the houses.}
            \end{itemize}
        \end{exbox}


    \section{Zeitliches \glqq vor\grqq}

        \subsection{Vor einem Zeitpunkt oder Ereignis}

            \begin{grammarentry}{before}{vor + Dativ}
                \GermanText{Zeitlich bezeichnet \textbf{vor} einen Zeitpunkt, der früher als ein anderer Zeitpunkt oder ein Ereignis liegt. In dieser Verwendung steht nach \textbf{vor} der Dativ.}
                \EnglishText{Temporally, \textbf{vor} refers to a point in time that is earlier than another point or event. In this use, \textbf{vor} takes the dative.}
            \end{grammarentry}

            \begin{exbox}
                \begin{itemize}
                    \ExampleItem{Vor dem Unterricht trinke ich Kaffee.}{I drink coffee before class.}
                    \ExampleItem{Bitte ruf mich vor der Besprechung an.}{Please call me before the meeting.}
                    \ExampleItem{Kurz vor Mitternacht ging er nach Hause.}{He went home shortly before midnight.}
                    \ExampleItem{Vor Beginn des Kurses müssen Sie sich anmelden.}{You must register before the course begins.}
                \end{itemize}
            \end{exbox}

        \subsection{Eine Zeitspanne zurück: \glqq ago\grqq}

            \begin{grammarentry}{ago}{vor + Zeitangabe}
                \GermanText{Mit einer Zeitspanne bezeichnet \textbf{vor}, wie lange ein vergangenes Ereignis zurückliegt. Die Zeitspanne steht im Dativ, aber ohne Artikel ist der Kasus an manchen Formen nicht sichtbar.}
                \EnglishText{With a period of time, \textbf{vor} indicates how long ago a past event occurred. The time expression is in the dative, although without an article the case is not visible in some forms.}

                \GermanText{Typische Struktur: \textbf{vor + Zahl + Zeiteinheit}.}
                \EnglishText{Typical structure: \textbf{vor + number + unit of time}.}
            \end{grammarentry}

            \begin{exbox}
                \begin{itemize}
                    \ExampleItem{Ich bin vor drei Tagen angekommen.}{I arrived three days ago.}
                    \ExampleItem{Sie hat vor einem Monat eine neue Arbeit begonnen.}{She started a new job one month ago.}
                    \ExampleItem{Vor vielen Jahren wohnte ich in dieser Stadt.}{Many years ago, I lived in this city.}
                \end{itemize}
            \end{exbox}

        \subsection{Uhrzeit}

            \begin{grammarentry}{fünf vor acht}{Angabe einer Uhrzeit}
                \GermanText{Bei Uhrzeiten bedeutet \textbf{vor}, dass eine bestimmte Zahl von Minuten bis zur genannten Stunde fehlt.}
                \EnglishText{With clock times, \textbf{vor} means that a certain number of minutes remain until the stated hour.}
            \end{grammarentry}

            \begin{exbox}
                \begin{itemize}
                    \ExampleItem{Es ist fünf vor acht.}{It is five minutes to eight.}
                    \ExampleItem{Der Zug fährt um Viertel vor neun ab.}{The train departs at a quarter to nine.}
                \end{itemize}
            \end{exbox}


    \section{\glqq vor\grqq und \glqq seit\grqq}

        \begin{grammarentry}{Vergangenheit}{ein Zeitpunkt oder eine Dauer}
            \GermanText{\textbf{vor} nennt den Abstand zu einem vergangenen Zeitpunkt. Die Handlung kann beendet sein.}
            \EnglishText{\textbf{vor} states the distance to a point in the past. The action may be finished.}

            \GermanText{\textbf{seit} bezeichnet den Beginn einer Situation, die bis zur Gegenwart andauert. Nach \textbf{seit} steht ebenfalls der Dativ.}
            \EnglishText{\textbf{seit} marks the beginning of a situation that continues up to the present. \textbf{seit} also takes the dative.}
        \end{grammarentry}

        \begin{exbox}
            \begin{itemize}
                \ExampleItem{Ich bin vor zwei Jahren nach Wien gezogen.}{I moved to Vienna two years ago.}
                \ExampleItem{Ich wohne seit zwei Jahren in Wien.}{I have been living in Vienna for two years.}
                \ExampleItem{Wir haben uns vor einer Woche getroffen.}{We met one week ago.}
                \ExampleItem{Wir kennen uns seit einer Woche.}{We have known each other for one week.}
            \end{itemize}
        \end{exbox}


    \section{Kausales \glqq vor\grqq}

        \begin{grammarentry}{Ursache einer Reaktion}{vor + Dativ}
            \GermanText{\textbf{vor} kann den Grund einer körperlichen, emotionalen oder unwillkürlichen Reaktion ausdrücken. Auch hier steht der Dativ.}
            \EnglishText{\textbf{vor} can express the cause of a physical, emotional, or involuntary reaction. The dative is also used here.}

            \GermanText{Häufige Verbindungen sind \textbf{vor Angst}, \textbf{vor Freude}, \textbf{vor Kälte}, \textbf{vor Hunger}, \textbf{vor Müdigkeit} und \textbf{vor Schmerz}.}
            \EnglishText{Common combinations include \textbf{vor Angst} (from fear), \textbf{vor Freude} (with joy), \textbf{vor Kälte} (from cold), \textbf{vor Hunger} (from hunger), \textbf{vor Müdigkeit} (from tiredness), and \textbf{vor Schmerz} (from pain).}
        \end{grammarentry}

        \begin{exbox}
            \begin{itemize}
                \ExampleItem{Das Kind weinte vor Angst.}{The child cried with fear.}
                \ExampleItem{Sie sprang vor Freude in die Luft.}{She jumped into the air with joy.}
                \ExampleItem{Meine Hände zittern vor Kälte.}{My hands are trembling from the cold.}
                \ExampleItem{Er konnte vor Müdigkeit kaum sprechen.}{He could hardly speak because of tiredness.}
            \end{itemize}
        \end{exbox}

        \begin{warningbox}
            \GermanText{Bei einer absichtlichen Begründung verwendet man normalerweise nicht \textbf{vor}, sondern zum Beispiel \textbf{wegen}, \textbf{aus} oder einen Nebensatz mit \textbf{weil}. \glqq Ich lerne vor der Prüfung\grqq{} bedeutet zeitlich \glqq before the exam\grqq, nicht \glqq because of the exam\grqq.}
            \EnglishText{For an intentional reason, German normally does not use \textbf{vor}, but for example \textbf{wegen}, \textbf{aus}, or a subordinate clause with \textbf{weil}. ``Ich lerne vor der Prüfung'' means temporally ``before the exam,'' not ``because of the exam.''}
        \end{warningbox}


    \section{Verben, Adjektive und Nomen mit \glqq vor\grqq}

        \begin{grammarentry}{Feste Rektion}{vor + Dativ}
            \GermanText{Manche Verben, Adjektive und Nomen verlangen eine Ergänzung mit \textbf{vor + Dativ}. In solchen Fällen muss die Verbindung als Einheit gelernt werden.}
            \EnglishText{Some verbs, adjectives, and nouns require a complement with \textbf{vor + dative}. In such cases, the combination should be learned as a unit.}
        \end{grammarentry}

        \rowcolors{2}{RowBlueLight}{RowBlueDark}
        \begin{longtable}{>{\bfseries}p{4cm}p{5.5cm}p{5.3cm}}
            \rowcolor{HeaderBlueDark}
            \color{white}\textbf{Verbindung} &
            \color{white}\textbf{Beispiel} &
            \color{white}\textbf{English} \\
            Angst haben vor & Ich habe Angst vor der Prüfung. & I am afraid of the exam. \\
            sich fürchten vor & Das Kind fürchtet sich vor dem Hund. & The child is afraid of the dog. \\
            warnen vor & Sie warnt uns vor der Gefahr. & She warns us about the danger. \\
            schützen vor & Die Jacke schützt mich vor der Kälte. & The jacket protects me from the cold. \\
            fliehen vor & Viele Menschen fliehen vor dem Krieg. & Many people flee from the war. \\
            sich verstecken vor & Er versteckt sich vor seinem Chef. & He is hiding from his boss. \\
            sicher vor & Die Daten sind vor unbefugtem Zugriff sicher. & The data is safe from unauthorized access. \\
            Respekt vor & Ich habe großen Respekt vor ihrer Arbeit. & I have great respect for her work. \\
            Vorrang vor & Sicherheit hat Vorrang vor Geschwindigkeit. & Safety takes priority over speed. \\
        \end{longtable}


    \section{Personen, Sachen und Pronominaladverbien}

        \begin{grammarentry}{vor ihm oder davor}{Person versus Sache}
            \GermanText{Bei Personen verwendet man \textbf{vor + Personalpronomen}: \textbf{vor mir}, \textbf{vor dir}, \textbf{vor ihm}, \textbf{vor ihr}, \textbf{vor uns}, \textbf{vor euch}, \textbf{vor ihnen}.}
            \EnglishText{With people, German uses \textbf{vor + personal pronoun}: \textbf{vor mir}, \textbf{vor dir}, \textbf{vor ihm}, \textbf{vor ihr}, \textbf{vor uns}, \textbf{vor euch}, \textbf{vor ihnen}.}

            \GermanText{Bei Sachen, Situationen oder ganzen Aussagen verwendet man häufig das Pronominaladverb \textbf{davor}. Die Frageform ist \textbf{wovor?}. Da \textbf{vor} mit einem Konsonanten beginnt, heißt es \textbf{davor} und \textbf{wovor}, nicht \glqq darvor\grqq{} oder \glqq worvor\grqq.}
            \EnglishText{With things, situations, or entire statements, German often uses the pronominal adverb \textbf{davor}. The question form is \textbf{wovor?}. Because \textbf{vor} begins with a consonant, the forms are \textbf{davor} and \textbf{wovor}, not ``darvor'' or ``worvor.''}
        \end{grammarentry}

        \begin{exbox}
            \begin{itemize}
                \ExampleItem{Ich habe Angst vor meinem Nachbarn. Ich habe Angst vor ihm.}{I am afraid of my neighbor. I am afraid of him.}
                \ExampleItem{Ich habe Angst vor der Prüfung. Ich habe Angst davor.}{I am afraid of the exam. I am afraid of it.}
                \ExampleItem{Wovor hast du Angst?}{What are you afraid of?}
                \ExampleItem{Ich warne dich davor, allein dorthin zu gehen.}{I warn you against going there alone.}
                \ExampleItem{Sie hat Angst davor, dass sie die Prüfung nicht besteht.}{She is afraid that she will not pass the exam.}
            \end{itemize}
        \end{exbox}

        \begin{grammarentry}{davor, dass / davor, zu}{Nebensatz und Infinitivgruppe}
            \GermanText{Eine Ergänzung kann durch \textbf{davor, dass + Nebensatz} ausgedrückt werden. Wenn das Subjekt beider Handlungen gleich ist, ist oft auch \textbf{davor, zu + Infinitiv} möglich.}
            \EnglishText{A complement can be expressed with \textbf{davor, dass + subordinate clause}. When the subject of both actions is the same, \textbf{davor, zu + infinitive} is often also possible.}
        \end{grammarentry}

        \begin{exbox}
            \begin{itemize}
                \ExampleItem{Er fürchtet sich davor, dass er seinen Arbeitsplatz verliert.}{He is afraid that he will lose his job.}
                \ExampleItem{Er fürchtet sich davor, seinen Arbeitsplatz zu verlieren.}{He is afraid of losing his job.}
            \end{itemize}
        \end{exbox}


    \section{\glqq vor\grqq, \glqq bevor\grqq und \glqq vorher\grqq}

        \begin{grammarentry}{Drei verschiedene Wortarten}{Präposition, Konjunktion, Adverb}
            \GermanText{\textbf{vor} ist hier eine Präposition und steht vor einem Nomen oder Pronomen: \textbf{vor dem Essen}.}
            \EnglishText{Here, \textbf{vor} is a preposition and stands before a noun or pronoun: \textbf{vor dem Essen} (before the meal).}

            \GermanText{\textbf{bevor} ist eine unterordnende Konjunktion. Sie leitet einen Nebensatz ein; das konjugierte Verb steht am Ende: \textbf{bevor ich esse}.}
            \EnglishText{\textbf{bevor} is a subordinating conjunction. It introduces a subordinate clause; the conjugated verb stands at the end: \textbf{bevor ich esse} (before I eat).}

            \GermanText{\textbf{vorher} ist ein Adverb und kann ohne nachfolgendes Nomen stehen: \textbf{Ich esse vorher}.}
            \EnglishText{\textbf{vorher} is an adverb and can stand without a following noun: \textbf{Ich esse vorher} (I eat beforehand).}
        \end{grammarentry}

        \rowcolors{2}{RowBlueLight}{RowBlueDark}
        \begin{longtable}{>{\bfseries}p{2.6cm}p{3.6cm}p{4.7cm}p{4.1cm}}
            \rowcolor{HeaderBlueDark}
            \color{white}\textbf{Wort} &
            \color{white}\textbf{Wortart} &
            \color{white}\textbf{Beispiel} &
            \color{white}\textbf{English} \\
            vor & Präposition & Vor dem Schlafengehen lese ich. & I read before going to sleep. \\
            bevor & Konjunktion & Bevor ich schlafen gehe, lese ich. & Before I go to sleep, I read. \\
            vorher & Adverb & Ich lese vorher. & I read beforehand. \\
        \end{longtable}


    \section{Wortstellung mit einer \glqq vor\grqq-Gruppe}

        \begin{grammarentry}{Satzglied}{Position im Hauptsatz}
            \GermanText{Eine Präpositionalgruppe mit \textbf{vor} ist ein Satzglied und kann oft verschoben werden. Steht sie im Vorfeld, folgt im deutschen Hauptsatz direkt das konjugierte Verb. Das Subjekt steht dann normalerweise nach dem Verb.}
            \EnglishText{A prepositional phrase with \textbf{vor} is a sentence element and can often be moved. When it occupies the first position, the conjugated verb follows directly in a German main clause. The subject then normally follows the verb.}
        \end{grammarentry}

        \begin{exbox}
            \begin{itemize}
                \ExampleItem{Ein Auto steht vor dem Haus.}{A car is standing in front of the house.}
                \ExampleItem{Vor dem Haus steht ein Auto.}{In front of the house, there is a car.}
                \ExampleItem{Ich trinke vor der Arbeit einen Kaffee.}{I drink a coffee before work.}
                \ExampleItem{Vor der Arbeit trinke ich einen Kaffee.}{Before work, I drink a coffee.}
            \end{itemize}
        \end{exbox}

        \begin{warningbox}
            \GermanText{Falsch ist: \glqq Vor der Arbeit ich trinke einen Kaffee.\grqq{} Nach dem ersten Satzglied muss im Hauptsatz das konjugierte Verb an zweiter Stelle stehen.}
            \EnglishText{Incorrect: ``Vor der Arbeit ich trinke einen Kaffee.'' After the first sentence element, the conjugated verb must be in second position in a main clause.}
        \end{warningbox}


    \section{Zusammenziehungen: \glqq vorm\grqq und \glqq vors\grqq}

        \begin{grammarentry}{Kurzformen}{vor dem / vor das}
            \GermanText{\textbf{vorm} ist die Zusammenziehung von \textbf{vor dem}. Sie ist besonders in der gesprochenen Sprache häufig und kommt auch in normaler geschriebener Sprache vor.}
            \EnglishText{\textbf{vorm} is the contraction of \textbf{vor dem}. It is especially common in spoken language and also occurs in ordinary written language.}

            \GermanText{\textbf{vors} ist die Zusammenziehung von \textbf{vor das}. Sie wird ebenfalls verwendet, wirkt aber stärker umgangssprachlich. In formellen Texten kann man die vollständige Form schreiben.}
            \EnglishText{\textbf{vors} is the contraction of \textbf{vor das}. It is also used, but sounds more colloquial. In formal texts, the full form can be written.}
        \end{grammarentry}

        \begin{exbox}
            \begin{itemize}
                \ExampleItem{Wir treffen uns vorm Kino.}{We are meeting in front of the cinema.}
                \ExampleItem{Stell den Stuhl vors Fenster.}{Put the chair in front of the window.}
            \end{itemize}
        \end{exbox}


    \section{Feste Ausdrücke mit \glqq vor\grqq}

        \begin{grammarentry}{Häufige Wendungen}{Bedeutung aus dem Zusammenhang}
            \GermanText{Einige häufige Wendungen müssen als feste Einheiten gelernt werden. Ihre Bedeutung lässt sich nicht immer nur aus der räumlichen Grundbedeutung erklären.}
            \EnglishText{Some common expressions should be learned as fixed units. Their meaning cannot always be explained solely from the basic spatial meaning.}
        \end{grammarentry}

        \rowcolors{2}{RowBlueLight}{RowBlueDark}
        \begin{longtable}{>{\bfseries}p{3.8cm}p{5.3cm}p{5.7cm}}
            \rowcolor{HeaderBlueDark}
            \color{white}\textbf{Ausdruck} &
            \color{white}\textbf{Bedeutung / Beispiel} &
            \color{white}\textbf{English} \\
            vor allem & hauptsächlich: Vor allem brauche ich Zeit. & above all; mainly: Above all, I need time. \\
            nach wie vor & weiterhin, noch immer: Das Problem besteht nach wie vor. & still; as before: The problem still exists. \\
            vor Ort & am betreffenden Ort: Ein Techniker ist vor Ort. & on site: A technician is on site. \\
            vor sich hin & ohne bestimmtes Ziel oder ohne andere zu beachten: Er redet vor sich hin. & to oneself / aimlessly: He is talking to himself. \\
            etwas vor Augen haben & sich etwas deutlich vorstellen: Ich habe das Ziel vor Augen. & to have something clearly in mind. \\
            jemandem etwas vor Augen führen & etwas deutlich zeigen oder erklären & to make something clear to someone. \\
            vor Kurzem & vor kurzer Zeit: Ich habe ihn vor Kurzem gesehen. & recently: I saw him recently. \\
        \end{longtable}


    \section{\glqq vor\grqq als Verbpräfix}

        \begin{grammarentry}{vor-}{meist betont und trennbar}
            \GermanText{Bei vielen Verben ist \textbf{vor-} ein betontes, trennbares Präfix. Im Hauptsatz steht der abgetrennte Präfixteil am Satzende. Im Infinitiv und im Nebensatz bleibt das Verb zusammengeschrieben. Im Partizip II steht häufig \textbf{ge} zwischen Präfix und Verbstamm.}
            \EnglishText{In many verbs, \textbf{vor-} is a stressed, separable prefix. In a main clause, the separated prefix appears at the end of the sentence. In the infinitive and in a subordinate clause, the verb remains written as one word. In the past participle, \textbf{ge} often stands between the prefix and the verb stem.}
        \end{grammarentry}

        \rowcolors{2}{RowBlueLight}{RowBlueDark}
        \begin{longtable}{>{\bfseries}p{3.1cm}p{4.5cm}p{4.5cm}p{3.1cm}}
            \rowcolor{HeaderBlueDark}
            \color{white}\textbf{Verb} &
            \color{white}\textbf{Hauptsatz} &
            \color{white}\textbf{Nebensatz} &
            \color{white}\textbf{Partizip II} \\
            vorlesen & Ich lese den Text vor. & ..., weil ich den Text vorlese. & vorgelesen \\
            vorbereiten & Ich bereite alles vor. & ..., weil ich alles vorbereite. & vorbereitet \\
            vorstellen & Ich stelle mich vor. & ..., weil ich mich vorstelle. & vorgestellt \\
            vorhaben & Ich habe heute viel vor. & ..., weil ich heute viel vorhabe. & vorgehabt \\
            vorkommen & Das kommt oft vor. & ..., weil das oft vorkommt. & vorgekommen \\
        \end{longtable}

        \begin{warningbox}
            \GermanText{Nicht jedes längere Verb, das mit den Buchstaben \textbf{vor} beginnt, folgt automatisch genau demselben Muster. Bei neuen Verben sollte man deshalb im Wörterbuch prüfen, ob das Verb trennbar ist und wie sein Partizip II lautet.}
            \EnglishText{Not every longer verb beginning with the letters \textbf{vor} automatically follows exactly the same pattern. Therefore, with new verbs, one should check in a dictionary whether the verb is separable and how its past participle is formed.}
        \end{warningbox}



    \section{Zeitgrenzen: vor oder bis?}
        \begin{grammarentry}{vor acht Uhr oder bis acht Uhr}{before versus by or until}
            \GermanText{\textbf{vor acht Uhr} bezeichnet einen Zeitpunkt, der früher als acht Uhr liegt. \textbf{bis acht Uhr} bezeichnet eine Frist oder das Ende einer Dauer; acht Uhr kann dabei dazugehören. \textbf{spätestens um acht Uhr} nennt den spätesten erlaubten Zeitpunkt.}
            \EnglishText{\textbf{vor acht Uhr} means earlier than eight o'clock. \textbf{bis acht Uhr} gives a deadline or an end point; eight o'clock may be included. \textbf{spätestens um acht Uhr} gives the latest allowed time.}
        \end{grammarentry}
        \begin{exbox}
            \begin{itemize}
                \ExampleItem{Bitte komm vor acht Uhr.}{Please come before eight o'clock.}
                \ExampleItem{Bitte komm bis acht Uhr.}{Please arrive by eight o'clock.}
                \ExampleItem{Ich arbeite bis acht Uhr.}{I work until eight o'clock.}
            \end{itemize}
        \end{exbox}

    \section{Vor einem substantivierten Verb}
        \begin{grammarentry}{vor dem Essen und vor dem Losfahren}{nominalised infinitives}
            \GermanText{Eine Handlung kann nach \textbf{vor} als substantivierter Infinitiv stehen: \textbf{vor dem Essen}, \textbf{vor dem Schlafengehen}, \textbf{vor dem Losfahren}. Der Infinitiv wird großgeschrieben und wie ein Nomen gebraucht. Für einen vollständigen Nebensatz verwendet man \textbf{bevor}.}
            \EnglishText{An action can follow \textbf{vor} as a nominalised infinitive: \textbf{vor dem Essen}, \textbf{vor dem Schlafengehen}, \textbf{vor dem Losfahren}. The infinitive is capitalised and used as a noun. Use \textbf{bevor} for a subordinate clause.}
        \end{grammarentry}
        \begin{exbox}
            \begin{itemize}
                \ExampleItem{Vor dem Losfahren kontrolliere ich die Reifen.}{Before setting off, I check the tyres.}
                \ExampleItem{Bevor ich losfahre, kontrolliere ich die Reifen.}{Before I set off, I check the tyres.}
                \ExampleItem{Vor dem Schlafengehen lese ich ein paar Seiten.}{Before going to sleep, I read a few pages.}
            \end{itemize}
        \end{exbox}

    \section{vor, vorher, zuvor und her in der Vergangenheit}
        \begin{grammarentry}{Zeitabstand oder Bezugspunkt}{ago versus earlier}
            \GermanText{\textbf{vor drei Tagen} zählt normalerweise vom Sprechzeitpunkt zurück. \textbf{drei Tage vorher} oder \textbf{drei Tage zuvor} bezieht sich dagegen auf einen bereits genannten Zeitpunkt. \textbf{Es ist drei Tage her, dass \ldots{}} kann ungefähr dieselbe Zeitspanne wie \textbf{vor drei Tagen} ausdrücken.}
            \EnglishText{\textbf{vor drei Tagen} normally counts back from the time of speaking. \textbf{drei Tage vorher} or \textbf{drei Tage zuvor} refers to an earlier reference point. \textbf{Es ist drei Tage her, dass \ldots{}} can express approximately the same interval as \textbf{vor drei Tagen}.}
        \end{grammarentry}
        \begin{exbox}
            \begin{itemize}
                \ExampleItem{Ich habe ihn vor drei Tagen gesehen.}{I saw him three days ago.}
                \ExampleItem{Am Freitag war die Prüfung. Drei Tage vorher habe ich viel gelernt.}{The exam was on Friday. Three days earlier, I studied a lot.}
                \ExampleItem{Es ist drei Tage her, dass ich ihn gesehen habe.}{It has been three days since I saw him.}
            \end{itemize}
        \end{exbox}
    \section{Typische Fehler}

        \begin{grammarentry}{Fehler 1}{falscher Kasus bei einer Position}
            \GermanText{Falsch: \glqq Das Auto steht vor das Haus.\grqq{} Richtig: \glqq Das Auto steht vor dem Haus.\grqq{} Das Auto befindet sich an einer Position; deshalb steht der Dativ.}
            \EnglishText{Incorrect: ``Das Auto steht vor das Haus.'' Correct: ``Das Auto steht vor dem Haus.'' The car is in a position; therefore, the dative is used.}
        \end{grammarentry}

        \begin{grammarentry}{Fehler 2}{falscher Kasus bei einer Platzierung}
            \GermanText{Falsch: \glqq Ich stelle das Auto vor dem Haus.\grqq{} Richtig, wenn eine neue Position gemeint ist: \glqq Ich stelle das Auto vor das Haus.\grqq{} Die Handlung beantwortet die Frage \glqq Wohin?\grqq.}
            \EnglishText{Incorrect: ``Ich stelle das Auto vor dem Haus.'' Correct, when a new position is meant: ``Ich stelle das Auto vor das Haus.'' The action answers the question ``Where to?''}
        \end{grammarentry}

        \begin{grammarentry}{Fehler 3}{\glqq vor\grqq statt \glqq bevor\grqq}
            \GermanText{Falsch: \glqq Vor ich gehe, rufe ich dich an.\grqq{} Richtig: \glqq Bevor ich gehe, rufe ich dich an.\grqq{} Vor einem vollständigen Nebensatz braucht man die Konjunktion \textbf{bevor}.}
            \EnglishText{Incorrect: ``Vor ich gehe, rufe ich dich an.'' Correct: ``Bevor ich gehe, rufe ich dich an.'' Before a complete subordinate clause, the conjunction \textbf{bevor} is required.}
        \end{grammarentry}

        \begin{grammarentry}{Fehler 4}{\glqq vor\grqq und \glqq seit\grqq verwechselt}
            \GermanText{Falsch für eine noch andauernde Situation: \glqq Ich wohne vor zwei Jahren hier.\grqq{} Richtig: \glqq Ich wohne seit zwei Jahren hier.\grqq{} Oder für ein abgeschlossenes Ereignis: \glqq Ich bin vor zwei Jahren hierhergezogen.\grqq}
            \EnglishText{Incorrect for a situation that is still continuing: ``Ich wohne vor zwei Jahren hier.'' Correct: ``Ich wohne seit zwei Jahren hier.'' Or for a completed event: ``Ich bin vor zwei Jahren hierhergezogen.''}
        \end{grammarentry}

        \begin{grammarentry}{Fehler 5}{Personen durch \glqq davor\grqq ersetzt}
            \GermanText{Bei Personen sagt man normalerweise \textbf{vor ihm}, \textbf{vor ihr} oder \textbf{vor ihnen}, nicht \textbf{davor}.}
            \EnglishText{With people, one normally says \textbf{vor ihm}, \textbf{vor ihr}, or \textbf{vor ihnen}, not \textbf{davor}.}
        \end{grammarentry}


    \newpage
    \section{Entscheidungsschema}

        \begin{grammarentry}{So wählst du die richtige Form}{Schritt für Schritt}
            \GermanText{1. Hat \textbf{vor} eine räumliche Bedeutung? Frage dann, ob eine Position oder eine gerichtete Platzierung gemeint ist.}
            \EnglishText{1. Does \textbf{vor} have a spatial meaning? Then ask whether a position or a directed placement is meant.}

            \GermanText{2. Position, Frage \glqq Wo?\grqq: Dativ. Gerichtete Platzierung, Frage \glqq Wohin?\grqq: Akkusativ.}
            \EnglishText{2. Position, question ``Where?'': dative. Directed placement, question ``Where to?'': accusative.}

            \GermanText{3. Hat \textbf{vor} eine zeitliche, kausale oder fest regierte Bedeutung? Dann steht normalerweise der Dativ.}
            \EnglishText{3. Does \textbf{vor} have a temporal, causal, or lexically governed meaning? Then the dative is normally used.}

            \GermanText{4. Folgt ein vollständiger Nebensatz? Dann brauchst du häufig \textbf{bevor} oder eine Verbindung mit \textbf{davor, dass}.}
            \EnglishText{4. Is a complete subordinate clause following? Then you often need \textbf{bevor} or a construction with \textbf{davor, dass}.}

            \GermanText{5. Bezieht sich die Ergänzung auf eine Sache? Dann kann oft \textbf{davor} stehen. Bezieht sie sich auf eine Person? Dann verwende \textbf{vor + Personalpronomen}.}
            \EnglishText{5. Does the complement refer to a thing? Then \textbf{davor} can often be used. Does it refer to a person? Then use \textbf{vor + personal pronoun}.}
        \end{grammarentry}



    % Erweiterung C1-C2: fortgeschrittene Bedeutungen und Konstruktionen
    \section{Räumliche Perspektive und Bezugsrahmen (Spatial perspective and reference frame)}
        \begin{grammarentry}{vor als Wechselpräposition}{perspective and location}
            \GermanText{Bei räumlichem \textbf{vor} hängt der Kasus davon ab, ob eine Lage (\textbf{Dativ}) oder ein Ziel der Positionierung (\textbf{Akkusativ}) gemeint ist. Bewegung allein erzwingt keinen Akkusativ. Welche Seite eines Gegenstandes als vorne gilt, hängt außerdem von dessen üblicher Orientierung, der Perspektive des Betrachters und dem Kontext ab.}
            \EnglishText{With spatial \textbf{vor}, case depends on whether a location (\textbf{dative}) or the target of positioning (\textbf{accusative}) is intended. Movement alone does not require accusative. What counts as the front of an object also depends on its conventional orientation, the observer's perspective and context.}
        \end{grammarentry}
        \begin{exbox}
            \begin{itemize}
                \ExampleItem{Die Demonstrierenden versammeln sich vor dem Rathaus.}{The demonstrators gather in front of the town hall.}
                \ExampleItem{Die Demonstrierenden stellen sich vor das Rathaus.}{The demonstrators position themselves in front of the town hall.}
                \ExampleItem{Ein Hund läuft vor dem Haus auf und ab.}{A dog is pacing back and forth in front of the house.}
            \end{itemize}
        \end{exbox}
        \begin{grammarentry}{vor, vorn und vorne}{preposition versus local adverb}
            \GermanText{\textbf{vor} verbindet sich als Präposition mit einer Ergänzung: \textbf{vor der Tür}. Dagegen sind \textbf{vorn} und \textbf{vorne} Lokaladverbien und stehen ohne ein solches Bezugsnomen. \textbf{nach vorn} drückt eine Richtung aus.}
            \EnglishText{The preposition \textbf{vor} takes a complement: \textbf{vor der Tür} (in front of the door). \textbf{vorn} and \textbf{vorne} are locative adverbs and do not need such a noun. \textbf{nach vorn} expresses direction.}
        \end{grammarentry}
        \begin{exbox}
            \begin{itemize}
                \ExampleItem{Der Ausgang befindet sich vorne.}{The exit is at the front.}
                \ExampleItem{Bitte treten Sie einen Schritt nach vorn.}{Please step forward.}
                \ExampleItem{Bitte warten Sie vor der Tür.}{Please wait in front of the door.}
            \end{itemize}
        \end{exbox}

    \section{Besondere räumliche und institutionelle Wendungen (Spatial and institutional expressions)}
        \begin{grammarentry}{vor Gericht und vor Augen}{fixed spatial interpretations}
            \GermanText{In festen Wendungen bezeichnet \textbf{vor} häufig nicht nur eine physische Position, sondern eine institutionelle oder bildliche Situation. Auch hier können \textbf{wo?} und \textbf{wohin?} verschiedene Kasus erfordern: \textbf{vor Gericht stehen} gegenüber \textbf{vor Gericht gehen}.}
            \EnglishText{In set expressions, \textbf{vor} often represents an institutional or figurative situation rather than merely a physical position. The distinction between location and destination can still determine case: \textbf{vor Gericht stehen} versus \textbf{vor Gericht gehen}.}
        \end{grammarentry}
        \begin{exbox}
            \begin{itemize}
                \ExampleItem{Der Angeklagte steht morgen vor Gericht.}{The defendant will appear in court tomorrow.}
                \ExampleItem{Der Streit soll vor Gericht gehen.}{The dispute is to be taken to court.}
                \ExampleItem{Die Ereignisse spielten sich vor den Augen zahlreicher Zeugen ab.}{The events unfolded before the eyes of numerous witnesses.}
                \ExampleItem{Sie stellte die Folgen des Beschlusses anschaulich vor Augen.}{She vividly illustrated the consequences of the decision.}
            \end{itemize}
        \end{exbox}

    \section{Nominalstil und komplexe Zeitangaben (Nominal style and complex temporal expressions)}
        \begin{grammarentry}{vor Beginn, vor Inkrafttreten}{formal article omission}
            \GermanText{Vor abstrakten Ereignisnomen wird in der gehobenen Amts- und Fachsprache der Artikel häufig weggelassen: \textbf{vor Beginn der Sitzung}, \textbf{vor Abschluss der Prüfung}, \textbf{vor Inkrafttreten des Gesetzes}. Die Ergänzung steht trotzdem im \textbf{Dativ}; das nachfolgende Nomen ist oft ein Genitivattribut.}
            \EnglishText{Formal administrative and specialist language often omits the article with abstract event nouns: \textbf{vor Beginn der Sitzung}, \textbf{vor Abschluss der Prüfung}, \textbf{vor Inkrafttreten des Gesetzes}. The complement still takes \textbf{dative}; the following noun is often a genitive modifier.}
        \end{grammarentry}
        \begin{exbox}
            \begin{itemize}
                \ExampleItem{Vor Inkrafttreten des Gesetzes sind Übergangsregelungen festzulegen.}{Transitional provisions must be defined before the law takes effect.}
                \ExampleItem{Vor Abschluss des Verfahrens dürfen keine Einzelheiten veröffentlicht werden.}{No details may be released before the procedure is concluded.}
            \end{itemize}
        \end{exbox}
        \begin{grammarentry}{vor zwei Jahren und von vor zwei Jahren}{reference time and attribution}
            \GermanText{\textbf{vor zwei Jahren} lokalisiert ein Ereignis relativ zur Sprechzeit in der Vergangenheit. In \textbf{die Unterlagen von vor zwei Jahren} gehört die Zeitangabe zu einer nominalen Charakterisierung. \textbf{zwei Jahre zuvor} bezieht sich dagegen üblicherweise auf einen bereits eingeführten Zeitpunkt der Erzählung.}
            \EnglishText{\textbf{vor zwei Jahren} locates an event relative to the time of speaking. In \textbf{die Unterlagen von vor zwei Jahren}, the time expression characterises a noun. \textbf{zwei Jahre zuvor}, by contrast, normally refers to a reference point already established in a narrative.}
        \end{grammarentry}
        \begin{exbox}
            \begin{itemize}
                \ExampleItem{Die Unterlagen von vor zwei Jahren sind noch verfügbar.}{The documents from two years ago are still available.}
                \ExampleItem{2019 wurde der Betrieb geschlossen; zwei Jahre zuvor war er erweitert worden.}{The company closed in 2019; two years earlier it had been expanded.}
            \end{itemize}
        \end{exbox}

    \section{Kausale Nuancen: vor, aus und wegen (Causal nuances: vor, aus and wegen)}
        \begin{grammarentry}{Unwillkürliche Reaktion und Beweggrund}{reaction versus motivation}
            \GermanText{Kausales \textbf{vor + Dativ} bezeichnet besonders häufig die Ursache einer körperlichen oder unwillkürlichen Reaktion: \textbf{vor Angst zittern}, \textbf{vor Freude weinen}. \textbf{aus + Dativ} benennt dagegen häufig einen Beweggrund für bewusstes Handeln: \textbf{aus Angst schweigen}. \textbf{wegen} ist als Begründung allgemeiner und hat andere Kasusregeln.}
            \EnglishText{Causal \textbf{vor + dative} very often expresses the cause of a physical or involuntary reaction: shaking with fear, crying with joy. \textbf{aus + dative} often gives a motive for intentional behaviour: remaining silent out of fear. \textbf{wegen} gives a more general reason and follows different case patterns.}
        \end{grammarentry}
        \begin{exbox}
            \begin{itemize}
                \ExampleItem{Sie zitterte vor Angst.}{She was trembling with fear.}
                \ExampleItem{Sie schwieg aus Angst vor Konsequenzen.}{She remained silent for fear of consequences.}
                \ExampleItem{Wegen des Sturms wurde die Veranstaltung abgesagt.}{The event was cancelled because of the storm.}
            \end{itemize}
        \end{exbox}
        \begin{grammarentry}{vor lauter, vor Neid, vor Begeisterung}{intensity and idiom}
            \GermanText{Die Verbindung \textbf{vor lauter + Nomen} verstärkt häufig eine unwillkürliche oder übermäßige Ursache. \textbf{vor lauter Aufregung} bedeutet sinngemäß wegen sehr großer Aufregung. Auch Wendungen wie \textbf{vor Neid erblassen} oder \textbf{vor Begeisterung außer sich sein} sind typisch.}
            \EnglishText{\textbf{vor lauter + noun} often emphasises an overwhelming cause. \textbf{vor lauter Aufregung} means roughly because of great excitement. Expressions such as \textbf{vor Neid erblassen} and \textbf{vor Begeisterung außer sich sein} are also common.}
        \end{grammarentry}
        \begin{exbox}
            \begin{itemize}
                \ExampleItem{Vor lauter Aufregung vergaß er seinen Text.}{He forgot his lines because he was so nervous.}
                \ExampleItem{Sie war vor Begeisterung ganz außer sich.}{She was beside herself with enthusiasm.}
            \end{itemize}
        \end{exbox}

    \section{Erweiterte Rektion und Pronominaladverbien (Advanced government and pronominal adverbs)}
        \begin{grammarentry}{Valenz mit vor und Dativ}{lexical government}
            \GermanText{Nicht jede Verbindung mit \textbf{vor} lässt sich rein räumlich oder zeitlich verstehen. Verb, Adjektiv oder Nomen können \textbf{vor + Dativ} als feste Ergänzung verlangen. Dazu gehören \textbf{sich hüten vor}, \textbf{jemanden warnen vor}, \textbf{jemanden schützen vor}, \textbf{Respekt vor} und \textbf{Vorrang vor}. Solche Rektionen müssen für jedes Wort gelernt werden.}
            \EnglishText{Not every \textbf{vor} phrase is purely spatial or temporal. A verb, adjective or noun may govern \textbf{vor + dative} as a fixed complement. Examples include \textbf{sich hüten vor}, \textbf{jemanden warnen vor}, \textbf{jemanden schützen vor}, \textbf{Respekt vor} and \textbf{Vorrang vor}. These patterns must be learnt word by word.}
        \end{grammarentry}
        \begin{exbox}
            \begin{itemize}
                \ExampleItem{Man sollte sich vor vorschnellen Schlussfolgerungen hüten.}{One should beware of hasty conclusions.}
                \ExampleItem{Die Fachleute warnten vor einer Unterschätzung des Risikos.}{The experts warned against underestimating the risk.}
                \ExampleItem{Der Datenschutz hat Vorrang vor bloßer Bequemlichkeit.}{Data protection takes precedence over mere convenience.}
            \end{itemize}
        \end{exbox}
        \begin{grammarentry}{davor, wovor, vor wem}{things, clauses and people}
            \GermanText{Bei Sachen und Sachverhalten stehen häufig \textbf{davor} und \textbf{wovor}; bei Personen verwendet man normalerweise \textbf{vor ihm}, \textbf{vor ihr} oder die Frage \textbf{vor wem?}. Eine ganze Handlung kann mit \textbf{davor, etwas zu tun} oder \textbf{davor, dass \ldots{}} angeschlossen werden.}
            \EnglishText{\textbf{davor} and \textbf{wovor} often refer to things or circumstances; people are usually expressed with \textbf{vor ihm}, \textbf{vor ihr} or the question \textbf{vor wem?}. A whole action can be added through \textbf{davor, etwas zu tun} or \textbf{davor, dass \ldots{}}.}
        \end{grammarentry}
        \begin{exbox}
            \begin{itemize}
                \ExampleItem{Wovor fürchtet sie sich? Vor dem Ergebnis.}{What is she afraid of? The result.}
                \ExampleItem{Vor wem fürchtet sie sich? Vor ihrem Vorgesetzten.}{Whom is she afraid of? Her superior.}
                \ExampleItem{Er warnte davor, die Ergebnisse vorschnell zu verallgemeinern.}{He warned against generalising the findings too hastily.}
            \end{itemize}
        \end{exbox}

    \section{Redewendungen, Präfix und Bedeutungswandel (Idioms, prefix and semantic extension)}
        \begin{grammarentry}{feste Ausdrücke}{idiomatic expressions}
            \GermanText{\textbf{vor allem} bedeutet insbesondere, \textbf{nach wie vor} weiterhin, \textbf{vor Ort} am betreffenden Ort, \textbf{von vornherein} von Anfang an, \textbf{vor sich hin} ohne bestimmten Adressaten oder Ziel. Diese Fügungen unterscheiden sich in Wortart und Satzbau von gewöhnlichen Präpositionalgruppen.}
            \EnglishText{\textbf{vor allem} means above all, \textbf{nach wie vor} still, \textbf{vor Ort} on site, \textbf{von vornherein} from the outset, and \textbf{vor sich hin} to oneself or aimlessly. Such expressions differ syntactically from ordinary prepositional phrases.}
        \end{grammarentry}
        \begin{exbox}
            \begin{itemize}
                \ExampleItem{Von vornherein war klar, dass der Zeitplan unrealistisch war.}{It was clear from the outset that the schedule was unrealistic.}
                \ExampleItem{Nach wie vor bestehen erhebliche Zweifel an der These.}{There are still considerable doubts about the thesis.}
                \ExampleItem{Er murmelte unverständliche Worte vor sich hin.}{He muttered unintelligible words to himself.}
            \end{itemize}
        \end{exbox}
        \begin{grammarentry}{vor- als trennbares Verbpräfix}{separable prefix}
            \GermanText{Das Verbpräfix \textbf{vor-} ist von der Präposition \textbf{vor} zu unterscheiden. Es kann eine vorausliegende Handlung, ein Vorzeigen oder eine Vorwegnahme ausdrücken. In \textbf{vorwerfen}, \textbf{vornehmen}, \textbf{vorsehen} und \textbf{vorziehen} sind Bedeutung und Rektion jeweils eigens zu lernen.}
            \EnglishText{The separable verb prefix \textbf{vor-} must be distinguished from the preposition \textbf{vor}. It can refer to an action done in advance, presentation or anticipation. The meanings and government of \textbf{vorwerfen}, \textbf{vornehmen}, \textbf{vorsehen} and \textbf{vorziehen} have to be learnt individually.}
        \end{grammarentry}
        \begin{exbox}
            \begin{itemize}
                \ExampleItem{Sie wirft ihm Nachlässigkeit vor.}{She accuses him of negligence.}
                \ExampleItem{Ich nehme mir vor, den Text noch einmal zu prüfen.}{I intend to review the text once more.}
                \ExampleItem{Der Entwurf sieht eine Übergangsfrist vor.}{The draft provides for a transitional period.}
                \ExampleItem{Ich ziehe eine sachliche Diskussion dem Streit vor.}{I prefer a factual discussion to an argument.}
            \end{itemize}
        \end{exbox}
    \section{Zusammenfassung}

        \begin{grammarentry}{Merksätze}{kurz und wichtig}
            \GermanText{Räumlich: \textbf{Wo? + Dativ}; \textbf{Wohin? + Akkusativ}.}
            \EnglishText{Spatially: \textbf{Where? + dative}; \textbf{Where to? + accusative}.}

            \GermanText{Zeitlich: \textbf{vor + Dativ} = before oder ago.}
            \EnglishText{Temporally: \textbf{vor + dative} = before or ago.}

            \GermanText{Kausal und in festen Verbindungen: normalerweise \textbf{vor + Dativ}.}
            \EnglishText{Causally and in fixed combinations: normally \textbf{vor + dative}.}

            \GermanText{Vor einem Nomen: \textbf{vor}; vor einem Nebensatz: häufig \textbf{bevor}; allein als Zeitadverb: \textbf{vorher}.}
            \EnglishText{Before a noun: \textbf{vor}; before a subordinate clause: often \textbf{bevor}; alone as a temporal adverb: \textbf{vorher}.}

            \GermanText{Sache: häufig \textbf{davor}; Person: \textbf{vor ihm / ihr / ihnen}.}
            \EnglishText{Thing: often \textbf{davor}; person: \textbf{vor ihm / ihr / ihnen}.}
        \end{grammarentry}

\end{document}
